\chapter{Filtrations and Information}

In financial modeling, and stochastic processes in general, the concept of \textit{information} is crucial. We need a mathematical way to represent "what is known" at a given time $t$. This is the role of a filtration.

\section{Filtrations}

\begin{definition}[Filtration]
A \textbf{filtration} on a probability space $(\Omega, \mathcal{F}, \mathbb{P})$ is an increasing family of sub-$\sigma$-algebras $(\mathcal{F}_t)_{t \in I}$ of $\mathcal{F}$. That is:
\[
\mathcal{F}_s \subseteq \mathcal{F}_t \subseteq \mathcal{F}, \quad \forall s \le t, \quad s, t \in I.
\]
A probability space equipped with a filtration, $(\Omega, \mathcal{F}, (\mathcal{F}_t)_{t \in I}, \mathbb{P})$, is called a \textbf{filtered probability space}.
\end{definition}

\subsection{Interpretation}
Each $\sigma$-algebra $\mathcal{F}_t$ represents the collection of events whose occurrence (or non-occurrence) is known by time $t$. The condition $\mathcal{F}_s \subseteq \mathcal{F}_t$ for $s \le t$ reflects the fact that information \textit{accumulates} over time: we do not forget what we knew in the past.
\begin{itemize}
    \item If $A \in \mathcal{F}_t$, then at time $t$, an observer knows whether the event $A$ has happened or not.
    \item $\mathcal{F}_0$ represents the initial information (often trivial, containing only $\emptyset$ and $\Omega$, if the process hasn't started or no prior information is available).
    \item $\mathcal{F}_T$ (or $\mathcal{F}_{\infty}$) represents the total information available at the end of the timeline.
\end{itemize}

\subsection{Example: The Coin Toss Filtration}
Recall our coin toss experiment with $\Omega = \{H, T\}^T$.
\begin{itemize}
    \item Let $\mathcal{F}_0 = \{ \emptyset, \Omega \}$. At the start, we know nothing about the outcome.
    \item Let $\mathcal{F}_t = \sigma(X_1, \dots, X_t)$. This is the $\sigma$-algebra generated by the first $t$ tosses.
\end{itemize}
Information at time $t$ consists of knowing the precise sequence of Heads and Tails up to step $t$, but nothing about steps $t+1, \dots, T$.
For $T=3$ and $t=1$:
\[
\mathcal{F}_1 = \sigma(X_1) = \{ \emptyset, \{H \times \{H,T\}^2\}, \{T \times \{H,T\}^2\}, \Omega \}
\]
Knowing which set in $\mathcal{F}_1$ contains $\omega$ is equivalent to knowing the value of $X_1(\omega)$.

\section{Adapted and Predictable Processes}

The relationship between a stochastic process and a filtration is fundamental. It tells us whether the value of the process is "knowable" at a certain time.

\subsection{Adapted Processes}
\begin{definition}[Adapted Process]
A stochastic process $X = (X_t)_{t \in I}$ is \textbf{adapted} with respect to the filtration $\mathbb{F} = (\mathcal{F}_t)_{t \in I}$ if closely linked to the filtration. specifically:
\[
X_t \text{ is } \mathcal{F}_t\text{-measurable for all } t \in I.
\]
\end{definition}

\textbf{Interpretation}: An adapted process extends the idea of "non-anticipative" behavior. The value $X_t$ is revealed at time $t$. We do not need to look into the future (time $u > t$) to know $X_t$. Historic stock prices are an adapted process: today's price is known today.

\begin{proposition}
If $X$ is adapted, then $X_s$ is $\mathcal{F}_t$-measurable for all $s \le t$.
\end{proposition}
\begin{proof}
Since $X$ is adapted, $X_s$ is $\mathcal{F}_s$-measurable. By definition of filtration, if $s \le t$, then $\mathcal{F}_s \subseteq \mathcal{F}_t$. Therefore, any function measurable with respect to $\mathcal{F}_s$ is automatically measurable with respect to the larger $\sigma$-algebra $\mathcal{F}_t$.
\end{proof}

\subsection{Predictable Processes}
\begin{definition}[Predictable Process]
A stochastic process $X = (X_t)_{t \in I}$ is \textbf{predictable} with respect to the filtration $\mathbb{F} = (\mathcal{F}_t)_{t \in I}$ if:
\[
X_{t} \text{ is } \mathcal{F}_{t-1}\text{-measurable for all } t \in I, \, t \ge 1.
\]
(Typically defining $X_0$ to be $\mathcal{F}_0$-measurable or constant).
\end{definition}

\textbf{Interpretation}: A predictable process is determined by the information available strictly \textit{before} time $t$. Its value at time $t$ is already "locked in" or "decied upon" at time $t-1$.
\begin{itemize}
    \item \textbf{Example}: Trading strategies. If you decide at time $t-1$ to hold $\phi_t$ shares of a stock during the interval $(t-1, t]$, this decision must differ on information available at $t-1$. Thus, the number of shares held is a predictable process.
    \item Note: In discrete time, predictability is a disjoint concept from purely adapted. Every predictable process is adapted (since $\mathcal{F}_{t-1} \subseteq \mathcal{F}_t$), but not every adapted process is predictable (e.g., the stock price itself is adapted but generally not predictable).
\end{itemize}

\section{Generated Filtration}

Every stochastic process generates its own "natural" flow of information.

\begin{definition}[Generated Filtration]
The \textbf{natural filtration} (or generated filtration) of a process $X$, denoted by $\mathbb{F}^X = (\mathcal{F}^X_t)_{t \in I}$, is defined as:
\[
\mathcal{F}^X_t := \sigma(X_s : s \le t) = \sigma(X_0, X_1, \dots, X_t).
\]
\end{definition}

\begin{proposition}
The generated filtration $\mathbb{F}^X$ is the \textbf{smallest} filtration with respect to which $X$ is adapted.
\end{proposition}
\begin{proof}
1. \textbf{Adaptedness}: By definition, $X_t$ is measurable with respect to $\sigma(X_0, \dots, X_t) = \mathcal{F}^X_t$. Thus $X$ is $\mathbb{F}^X$-adapted.
2. \textbf{Minimality}: Let $\mathbb{G} = (\mathcal{G}_t)$ be any other filtration such that $X$ is $\mathbb{G}$-adapted. Then for any $t$, $X_0, \dots, X_t$ are all $\mathcal{G}_t$-measurable (since $X_s \in \mathcal{G}_s \subseteq \mathcal{G}_t$). Since $\mathcal{F}^X_t$ is the \textit{smallest} $\sigma$-algebra making $X_0, \dots, X_t$ measurable, it follows that $\mathcal{F}^X_t \subseteq \mathcal{G}_t$.
\end{proof}

Currently, we often work with the natural filtration unless specified otherwise (e.g., if there are multiple sources of randomness or "insider information").
